\section{Selecting a coefficient by a permanent projection}
\label{sec:selector}

We realize a coefficient of a diagonally shifted permanent as an affine
projection of a larger permanent.  The construction selects prescribed
exponents in two auxiliary variables while keeping the larger permanent's
size linear in the order of the shifted permanent.

Fix integers $k\ge1$ and $0\le q\le\lfloor k/2\rfloor$, and put
$r=k-2q$.  Let $\theta_1,\ldots,\theta_k$ be indeterminates and put
$D=\operatorname{diag}(\theta_1,\ldots,\theta_k)$.  Let
$X=(X_{ij})_{1\le i,j\le k}$ be a matrix of independent variables.
Define the integral coefficient polynomial
$G_D\in\mathbb Z[\theta_1,\ldots,\theta_k][X_{ij}:1\le i,j\le k]$ by
\begin{equation}
 G_D(X)=[s^qt^q]\per_k(X+sI+tD).
 \label{eq:selector-coefficient}
\end{equation}
The brackets denote coefficient extraction in the auxiliary variables
$s,t$.  For a tuple of values $\theta_i$ in any field, we use the same
notation $G_D$ for the specialized polynomial.  Write $[k]=\{1,\ldots,k\}$.
Expanding the diagonal entries gives
\begin{equation}
 G_D(X)=
 \sum_{\substack{S,T\subseteq[k],\ |S|=|T|=q\\S\cap T=\varnothing}}
 \left(\prod_{i\in T}\theta_i\right)\per(X_{C,C}),
 \qquad C=[k]\setminus(S\cup T).
 \label{eq:selector-subsets}
\end{equation}
Indeed, $S$ records the diagonal positions contributing $s$, and $T$
those contributing $t\theta_i$.  The two sets are disjoint because each
row and column is used once in a permanent term.  After these fixed
diagonal positions are removed, the remaining permutation is an arbitrary
permutation of $C$, and these data recover the original term uniquely.
Thus no additional factorial occurs in
\eqref{eq:selector-subsets}.  Every summand is homogeneous of degree $r$
in $X$.
We use the convention that the permanent of the empty matrix is $1$.

We use the matching interpretation of a permanent.  A matrix's rows and
columns form the two vertex classes of a bipartite graph, with each entry
as the weight of the edge joining its row to its column.  The permanent is the
sum, over perfect matchings, of the products of their edge weights.
Negative weights are allowed.  Weighted matching gadgets and cancellation are
classical tools of Valiant's projection framework \cite{Valiant1979Algebra}.
The practice of fixing external choices before summing internal
completions also appears in his permanent reduction
\cite[Lemma~3.1 and its proof]{Valiant1979Permanent}.
Signed gadgets enforcing equality of edge choices also appear in
Dell, Husfeldt, Marx, Taslaman and Wahl\'en
\cite[Section~3 and Figure~2]{DellHusfeldtMarxTaslamanWahlen2014}.
Here we give the weights and completion counts needed for a construction
of linear size.

\begin{proposition}[Coefficient projection]\label{prop:selector}
Let $k\ge1$ and $0\le q\le\lfloor k/2\rfloor$ be integers, put $r=k-2q$,
and let $D=\operatorname{diag}(\theta_1,\ldots,\theta_k)$ with
$\theta_i\in\C$.  For every integer $m\ge11k-2q$, there is an $m\times m$
matrix $A(X)$ of complex affine-linear forms such that
\begin{equation}
 \per_m(A(X))=a_0G_D(X),
 \qquad a_0=(-2)^{2k}\big((k-q)!\big)^4\ne0.
 \label{eq:selector-scalar}
\end{equation}
Its entrywise homogenization $B(X,z)$ is a matrix of homogeneous linear
forms whose permanent is the degree-$m$ homogenization of $a_0G_D$:
\begin{equation}
 \per_m(B(X,z))=a_0z^{m-r}G_D(X).
 \label{eq:selector-homogeneous}
\end{equation}
The construction requires no distinctness or nonvanishing assumption on
the entries of $D$.
\end{proposition}

\begin{proof}
Start with $k$ main rows and columns carrying the edges $X_{ij}$.
At each diagonal position add two labeled special edges, with weights
\[
 w_{i,1}=1,\qquad w_{i,2}=\theta_i.
\]
We call them the first and second color.  At this intermediate stage
they are parallel to the ordinary edge $X_{ii}$.  The three labels record
the choices $X_{ii}$, $s$, and $t\theta_i$ in the diagonal entry, with the
factors $s,t$ omitted from the two special weights.  The coefficient of
$s^qt^q$ in the shifted permanent is therefore the weighted sum of main
matchings using exactly $q$ special edges of each color, as in
\eqref{eq:selector-subsets}.
We will count each color in a separate gate and couple each special
main edge to the corresponding distinguished edge of its gate.

\emph{A gate imposing one prescribed count.}
Put $h=k-q$.  Take $k$ original rows and columns, with distinguished
diagonal edges of weight $1$, and $h$ additional rows and columns.
Join every original row to every additional column and every additional
row to every original column, all with weight $1$.  There are no other
edges.  Every additional row must occupy an original column, and every
additional column must be occupied by an original row.  Consequently
exactly $q$ distinguished edges occur in a perfect matching.

For each prescribed distinguished set $S\subseteq[k]$ of size $q$, the
remaining original rows match the additional columns in $h!$ ways, and
the additional rows match the remaining original columns in $h!$ ways.
These choices are independent.  Hence the multiplicity for this exact
set is
\begin{equation}
 (h!)^2=((k-q)!)^2.
 \label{eq:selector-gate-count}
\end{equation}
This also covers $q=0$: all original rows and columns then match across
the two groups.

\emph{A gadget equating two edge choices.}
Call the main and gate vertices \emph{outer}.  For one special main edge
of weight $w$ and its distinguished gate edge
of weight $1$, write the outer row endpoints as $u_1,u_2$ and the outer
column endpoints as $v_1,v_2$, respectively.  These two edges have distinct
endpoints.  Delete them and add three fresh internal rows and three fresh
internal columns, with weight matrix
\begin{equation}
 H=\begin{pmatrix}-1&-1&1\\-1&-1&1\\1&1&1\end{pmatrix}.
 \label{eq:selector-gadget}
\end{equation}
For $i=1,2$, add an incoming edge from $u_i$ to internal column $i$ and
an outgoing edge from internal row $i$ to $v_i$.  All four weights are
$1$, except that the first incoming edge has weight $-2w$.

Fix the subset $I\subseteq\{1,2\}$ of incoming edges used and the
subset $O\subseteq\{1,2\}$ of outgoing edges used.  Let $\Phi_w(I,O)$
be the total weight of their internal completions, including the weights
of those boundary edges.  If $|I|\ne|O|$,
the remaining internal row and column sets have different sizes, so
there is no completion and $\Phi_w(I,O)=0$.  Otherwise
\begin{equation}
 \Phi_w(I,O)=(-2w)^{\mathbf1_{\{1\in I\}}}
 \per H_{\{1,2,3\}\setminus O,\,\{1,2,3\}\setminus I}.
 \label{eq:selector-boundary-weight}
\end{equation}
The complete table of $\Phi_w$ is
\begin{equation}
 \begin{array}[t]{c|rrrr}
 I\backslash O&\varnothing&\{1\}&\{2\}&\{1,2\}\\\hline
 \varnothing&-2&0&0&0\\
 \{1\}&0&0&0&0\\
 \{2\}&0&0&0&0\\
 \{1,2\}&0&0&0&-2w.
 \end{array}
 \label{eq:selector-boundary-table}
\end{equation}
To verify the six balanced cases, first $\per H=-2$.  Deleting any one
row and any one column indexed in $\{1,2\}$ leaves
$\left(\begin{smallmatrix}-1&1\\1&1\end{smallmatrix}\right)$, whose
permanent is $-1+1=0$.  This includes the crossed choices
$I=\{1\},O=\{2\}$ and $I=\{2\},O=\{1\}$.  Deleting both interface
rows and columns leaves the entry $H_{33}=1$; the four external edge
weights then have product $-2w$.  The ten unbalanced cases have already
been excluded.

Thus $\Phi_w(I,O)$ vanishes unless $I=O=\varnothing$ or
$I=O=\{1,2\}$.  In these two states the weights $-2$ and $-2w$
reproduce neither outer edge being used or both being used, with the
common factor $-2$.  The four balanced singleton states can have
individual internal matchings; it is their sums that cancel.

\emph{Combining the two counts.}
Take two disjoint copies of the count gate, one for each color.
For every $i\in[k]$ and $c\in\{1,2\}$, replace the special main edge
of weight $w_{i,c}$ and its distinguished gate edge by the equality
gadget.  Use fresh internal rows and columns for all $2k$ gadgets.
At a fixed index $i$, the two replacements share the main row and column,
while their internal vertices and their gate endpoints remain distinct.
Figure~\ref{fig:selector-wiring} shows these attachments.

\begin{figure}[tbp]
\centering
\begingroup
\setlength{\unitlength}{1pt}
\begin{picture}(460,202)
\linethickness{0.45pt}
\put(174,129){\framebox(112,55){}}
\put(174,25){\framebox(112,55){}}
\put(72,104){\line(1,0){316}}
\qbezier(72,104)(123,122)(174,140)
\qbezier(286,140)(337,122)(388,104)
\put(72,171){\line(1,0){102}}
\put(286,171){\line(1,0){102}}
\qbezier(72,104)(123,86)(174,68)
\qbezier(286,68)(337,86)(388,104)
\put(72,37){\line(1,0){102}}
\put(286,37){\line(1,0){102}}
\put(72,104){\circle*{3.2}}
\put(388,104){\circle*{3.2}}
\put(72,171){\circle*{3.2}}
\put(388,171){\circle*{3.2}}
\put(72,37){\circle*{3.2}}
\put(388,37){\circle*{3.2}}
\put(174,171){\circle*{3.2}}
\put(174,140){\circle*{3.2}}
\put(286,171){\circle*{3.2}}
\put(286,140){\circle*{3.2}}
\put(174,68){\circle*{3.2}}
\put(174,37){\circle*{3.2}}
\put(286,68){\circle*{3.2}}
\put(286,37){\circle*{3.2}}
\put(230,162){\makebox(0,0){\small color 1: \(H\)}}
\put(230,151){\makebox(0,0){\scriptsize 3 rows \(+\) 3 columns}}
\put(181,171){\makebox(0,0)[l]{\scriptsize \(c_2\)}}
\put(181,140){\makebox(0,0)[l]{\scriptsize \(c_1\)}}
\put(279,171){\makebox(0,0)[r]{\scriptsize \(r_2\)}}
\put(279,140){\makebox(0,0)[r]{\scriptsize \(r_1\)}}
\put(230,57){\makebox(0,0){\small color 2: \(H\)}}
\put(230,46){\makebox(0,0){\scriptsize 3 rows \(+\) 3 columns}}
\put(181,68){\makebox(0,0)[l]{\scriptsize \(c_1\)}}
\put(181,37){\makebox(0,0)[l]{\scriptsize \(c_2\)}}
\put(279,68){\makebox(0,0)[r]{\scriptsize \(r_1\)}}
\put(279,37){\makebox(0,0)[r]{\scriptsize \(r_2\)}}
\put(64,104){\makebox(0,0)[r]{\shortstack[r]{\small shared main\\[-1pt]\small row \(i\)}}}
\put(396,104){\makebox(0,0)[l]{\shortstack[l]{\small shared main\\[-1pt]\small column \(i\)}}}
\put(72,192){\makebox(0,0){\small gate 1 row \(i\)}}
\put(388,192){\makebox(0,0){\small gate 1 column \(i\)}}
\put(72,16){\makebox(0,0){\small gate 2 row \(i\)}}
\put(388,16){\makebox(0,0){\small gate 2 column \(i\)}}
\put(230,113){\makebox(0,0){\small \(X_{ii}\)}}
\put(122,132){\makebox(0,0){\small \(-2\)}}
\put(121,76){\makebox(0,0){\small \(-2\theta_i\)}}
\end{picture}
\endgroup
\caption{Wiring after the two replacements at a fixed diagonal position
\(i\). Each box is a separate copy of \(H\) on three fresh rows and
three fresh columns. Its incoming ports \(c_1,c_2\) are internal columns
and its outgoing ports \(r_1,r_2\) are internal rows; their placement is
schematic. The two copies share the main endpoints and have distinct
gate and internal vertices. Unlabeled boundary edges have weight \(1\).
The ordinary edge \(X_{ii}\) remains; other main and gate edges are omitted.}
\label{fig:selector-wiring}
\end{figure}


Group the full matching sum by the set $E$ of chosen edges outside the
internal $H$-blocks.
This set includes all chosen incoming and outgoing edges and must cover
each outer vertex exactly once.  Let $w_{\mathrm{out}}(E)$ be the product
of the weights of the chosen edges whose two endpoints are outer.
For the gadget indexed by $(i,c)$, write $I_{i,c},O_{i,c}$ for its
boundary state.  Since the internal vertex sets are disjoint, the total
weight of all completions of $E$ is
\[
 w_{\mathrm{out}}(E)
 \prod_{i=1}^k\prod_{c=1}^2
 \Phi_{w_{i,c}}(I_{i,c},O_{i,c}).
\]
Each factor $\Phi_{w_{i,c}}$ already includes its boundary edge weights;
$w_{\mathrm{out}}(E)$ contains only the remaining chosen weights.
By \eqref{eq:selector-boundary-table}, this
product vanishes unless every gadget has $I_{i,c}=O_{i,c}=\varnothing$
or $I_{i,c}=O_{i,c}=\{1,2\}$.  Unbalanced states have no completion,
and balanced singleton states cancel.

Now restrict to choices $E$ for which every gadget has one of these two states.
For each gadget with $I_{i,c}=O_{i,c}=\{1,2\}$, replace its four boundary
edges by the original special main edge and distinguished gate edge.
For an empty state, insert neither original edge.  Because $E$ covers every
outer vertex exactly once, this contraction gives a perfect matching in the
graph before replacement, with each special main edge used if and only if
its paired gate edge is used.  Conversely, each such matching determines
the set $E$ and its boundary states uniquely.  The table assigns $-2w_{i,c}$
to a used pair of original weight $w_{i,c}$, and $-2$ to an unused pair of weight
$1$.  Hence its contribution is its original matching weight multiplied
by exactly one factor $-2$ for each of the $2k$ gadgets.  Shared main
endpoints enter through the matching constraint enforced when $E$ is fixed;
the internal factorization uses disjoint vertices.

Each gate now forces exactly $q$ special edges of its color.  The two
selected sets are disjoint by the main matching condition.  For these
sets, the two gates contribute $((k-q)!)^4$ by
\eqref{eq:selector-gate-count}, while the main matching sum is the
corresponding summand of \eqref{eq:selector-subsets}.  The $2k$ gadgets
contribute $(-2)^{2k}$.  This proves \eqref{eq:selector-scalar} for the
unpadded construction.

Its number of rows, and separately its number of columns, is
\begin{equation}
 k+2(2k-q)+3(2k)=11k-2q.
 \label{eq:selector-size}
\end{equation}
These terms count the main graph, the two count gates, and the fresh
internal blocks.  The final graph has no parallel special edges: their
replacements lead to distinct internal columns.  Its weight matrix
therefore has entries that are variables or constants.  Append an
identity block to reach order $m$; its unique matching has weight $1$,
so the permanent identity is preserved.

Finally, write this padded affine matrix as
$A(X)=A_0+\sum_{i,j}X_{ij}A_{ij}$ and set
$B(X,z)=zA_0+\sum_{i,j}X_{ij}A_{ij}$.  For $z\ne0$, homogeneity gives
\[
 \per_m B(X,z)
 =z^m\per_m A(X/z)
 =a_0z^mG_D(X/z)
 =a_0z^{m-r}G_D(X).
\]
Since $m\ge k\ge r$, both sides are polynomials.  Their equality for
$z\ne0$ proves their identity also at $z=0$.  The identity block is
homogenized along with every other constant entry, so it contributes
the additional power of $z$ required in
\eqref{eq:selector-homogeneous}.
\end{proof}

The affine identity yields a comparison for every fixed affine substitution.

\begin{corollary}\label{cor:selector-border}
Under the hypotheses of Proposition~\ref{prop:selector}, let $d\ge1$ be an
integer and let
$\Lambda:\C^d\to\Mat_k(\C)$ be any fixed affine map.  Then
\[
 \bdc(G_D\circ\Lambda)\le\bdc(\per_m).
\]
In particular, this holds for every fixed linear restriction of $G_D$.
\end{corollary}

\begin{proof}
Proposition~\ref{prop:selector} gives
$G_D\circ\Lambda=a_0^{-1}\per_m\circ A\circ\Lambda$.
Apply Lemma~\ref{lem:affine-substitution} to the affine map
$A\circ\Lambda$ and the nonzero scalar $a_0^{-1}$.
Both are fixed before any determinant approximation is chosen.
\end{proof}

\begin{remark}\label{rem:selector-characteristic}
The coefficient formula \eqref{eq:selector-coefficient} and the subset
formula \eqref{eq:selector-subsets} are integral and can be specialized
over any field.  The projection used for the complex border comparison
requires the nonzero scalar $a_0$, whose factor $(-2)^{2k}$ vanishes in
characteristic two.  The characteristic-two argument uses the integral
coefficient expression specialized directly to $\overline{\mathbb F}_2$.
\end{remark}
